
\documentclass{article}
\usepackage{colortbl}
\usepackage{makecell}
\usepackage{multirow}
\usepackage{supertabular}

\begin{document}

\newcounter{utterance}

\centering \large Interaction Transcript for game `eqbench', experiment `eq\_bench', episode 131 with reflexion{-}llama{-}3.1{-}8b{-}instruct.
\vspace{24pt}

{ \footnotesize  \setcounter{utterance}{1}
\setlength{\tabcolsep}{0pt}
\begin{supertabular}{c@{$\;$}|p{.15\linewidth}@{}p{.15\linewidth}p{.15\linewidth}p{.15\linewidth}p{.15\linewidth}p{.15\linewidth}}
    \# & \multicolumn{2}{c}{Player} && \multicolumn{2}{c}{Game Master} \\
    \hline

    \theutterance \stepcounter{utterance}  
    & & & \multicolumn{4}{p{0.6\linewidth}}{
        \cellcolor[rgb]{0.9,0.9,0.9}{
            \makecell[{{p{\linewidth}}}]{
                \texttt{\tiny{[P1$\langle$GM]}}
                \texttt{You have attempted this same task before and did not succeed. These are the plans you wrote after those attempts.} \\
\texttt{Trial \#1: New plan: I failed to accurately predict Sam's emotional response because I didn't consider the context of the conversation and the history of their relationship. Specifically, I didn't take into account Sam's past experience with Arthur's "creative decisions" (painting the treehouse neon green) and how it might have affected their trust and communication. Next time, I will consider the following:} \\
\\ 
\texttt{1. Analyze the conversation history to identify any past conflicts or issues that might be relevant to the current situation.} \\
\texttt{2. Consider the character's personality, values, and motivations to better understand their emotional response.} \\
\texttt{3. Look for specific phrases or actions that might indicate a particular emotional response (e.g., Sam's mention of the treehouse and the wedding speech suggests that he is still hurt by Arthur's past actions).} \\
\\ 
\texttt{By taking these factors into account, I will be able to provide a more accurate prediction of Sam's emotional response.} \\
\\ 
\\ 
\\ 
\texttt{Your task is to predict the likely emotional responses of a character in this dialogue:} \\
\\ 
\texttt{Arthur: You remember that time we built that treehouse in the back garden, right?} \\
\texttt{Sam: What's that got to do with anything?} \\
\texttt{Arthur: I'm just saying, we were a team, you and me. We built something out of nothing. } \\
\texttt{Sam: Yeah, and then you decided it would be fun to paint it neon green without asking me. } \\
\texttt{Arthur: It was a creative decision, Sammy. } \\
\texttt{Sam: It was a stupid decision, Arthur. Just like your speech at my wedding.} \\
\texttt{Arthur: That's not fair, Sam. I messed up, I know. But I'm trying to apologize here.} \\
\texttt{Sam: And I'm supposed to just accept it because you've decided to say sorry? } \\
\texttt{Arthur: No, not because I've decided to. Because I mean it, Sam. I'm sorry.} \\
\texttt{{[}End dialogue{]}} \\
\\ 
\texttt{At the end of this dialogue, Sam would feel...} \\
\texttt{Pacified} \\
\texttt{Skeptical} \\
\texttt{Vindicated} \\
\texttt{Joyful} \\
\\ 
\texttt{Give each of these possible emotions a score from 0{-}10 for the relative intensity that they are likely to be feeling each.} \\
\\ 
\texttt{You must output in the following format, including headings (of course, you should give your own scores), with no additional commentary:} \\
\\ 
\texttt{Pacified: <score>} \\
\texttt{Skeptical: <score>} \\
\texttt{Vindicated: <score>} \\
\texttt{Joyful: <score>} \\
\\ 
\\ 
\texttt{{[}End of answer{]}} \\
\\ 
\texttt{Remember: zero is a valid score, meaning they are likely not feeling that emotion. You must score at least one emotion > 0.} \\
\\ 
\texttt{Your answer:} \\
            }
        }
    }
     \\ \\

    \theutterance \stepcounter{utterance}  
    & \multicolumn{4}{p{0.6\linewidth}}{
        \cellcolor[rgb]{0.561,0.737,0.561}{
            \makecell[{{p{\linewidth}}}]{
                \texttt{\tiny{[P1$\rangle$GM]}}
                \texttt{Pacified: 2} \\
\texttt{Skeptical: 8} \\
\texttt{Vindicated: 0} \\
\texttt{Joyful: 0} \\
            }
        }
    }
    & & \\ \\

    \theutterance \stepcounter{utterance}  
    & & & \multicolumn{2}{p{0.3\linewidth}}{
        \cellcolor[rgb]{0.9,0.9,0.9}{
            \makecell[{{p{\linewidth}}}]{
                \texttt{\tiny{[GM$|$GM]}}
                \texttt{Pacified: 2} \\
\texttt{Skeptical: 8} \\
\texttt{Vindicated: 0} \\
\texttt{Joyful: 0} \\
            }
        }
    }
    & & \\ \\

    \theutterance \stepcounter{utterance}  
    & & & \multicolumn{2}{p{0.3\linewidth}}{
        \cellcolor[rgb]{0.9,0.9,0.9}{
            \makecell[{{p{\linewidth}}}]{
                \texttt{\tiny{[GM$|$GM]}}
                \texttt{Pacified: 4} \\
\texttt{Skeptical: 7} \\
\texttt{Vindicated: 3} \\
\texttt{Joyful: 0} \\
            }
        }
    }
    & & \\ \\

    \theutterance \stepcounter{utterance}  
    & & & \multicolumn{2}{p{0.3\linewidth}}{
        \cellcolor[rgb]{0.9,0.9,0.9}{
            \makecell[{{p{\linewidth}}}]{
                \texttt{\tiny{[GM$|$GM]}}
                \texttt{game\_result = LOSE} \\
            }
        }
    }
    & & \\ \\

    \theutterance \stepcounter{utterance}  
    & & \multicolumn{4}{p{0.6\linewidth}}{
        \cellcolor[rgb]{0.95,0.95,0.95}{
            \makecell[{{p{\linewidth}}}]{
                \texttt{\tiny{[P1$\leftrightarrow$P]}}
                \texttt{Post{-}episode Reflexion (not sent to GM):} \\
\texttt{New plan: In my previous attempt, I scored Sam's emotional response as Pacified: 2, but this was likely an underestimation. Upon re{-}examining the dialogue, I realize that Sam's skepticism is even more intense due to the mention of the wedding speech, which suggests a deeper hurt and betrayal. I should have also considered the fact that Sam's response to Arthur's apology is not just about accepting or rejecting it, but also about the sincerity and effort Arthur puts into making amends.} \\
\\ 
\texttt{To improve my prediction, I will pay closer attention to the following:} \\
\\ 
\texttt{1. The tone and language used by the character in their response, which can indicate the intensity of their emotions.} \\
\texttt{2. The specific events or actions mentioned in the conversation that might be relevant to the character's emotional response.} \\
\texttt{3. The character's past behavior and interactions with others, which can provide insight into their emotional patterns and reactions.} \\
\\ 
\texttt{By taking these factors into account, I will be able to provide a more accurate prediction of Sam's emotional response, likely scoring Skeptical: 9 or 10.} \\
            }
        }
    }
    & \\ \\

\end{supertabular}
}

\end{document}
